Benchmark contamination---the presence of test data in training corpora---is extensively documented in large language model evaluation, yet no prior work quantifies how contamination translates to score inflation. We address this gap through checkpoint-gradient analysis: by comparing model checkpoints at different training stages with varying contamination exposure, we measure the correlation between contamination and benchmark performance without requiring contamination-free baselines. Using capability detrending via out-of-distribution perplexity, we isolate contamination effects from legitimate capability gains. In preliminary validation across 72 Pythia checkpoint-benchmark pairs (6 model sizes $\times$ 12 checkpoints), our methodology yields a positive correlation (Spearman $r = 0.326$, $p = 0.003$) between contamination exposure and inflation residual. Per-benchmark analysis suggests stronger effects for factual benchmarks (MMLU: $r = 0.41$, $p < 0.01$) compared to format-resistant tasks (WinoGrande: $r = 0.18$, $p = 0.15$, not significant). Our n-gram overlap analysis confirms detection mechanism validity, with individual MMLU items showing up to 17.4\% overlap with The Pile training corpus. These preliminary findings suggest that contamination-inflation correlation is measurable, opening possibilities for contamination-adjusted benchmark evaluation pending full-scale validation.
